% Terjemahan Bahasa Indonesia dari chapter1.tex baris 33-139.
% Sumber: Methods of Algebra, Volume 2, commit
% 9a5803ff77dd3257484cb177f851a73770a59dd3 (CC BY 4.0).
% stable-unit-id: o014.aljabr2.chapter1.subquotients

% segment-id: o014.li2.chapter1.subquotients.h001
\section{Subkuosien}\label{sec:subquot}
\hypertarget{o014.aljabr2.chapter1.subquotients}{ }

% segment-id: o014.li2.chapter1.subquotients.p001
Pertama-tama, kita mengingat kembali pengertian monomorfisme dan
epimorfisme. Tetapkan sebuah kategori $\mathcal{C}$.

% segment-id: o014.li2.chapter1.subquotients.l001
\begin{itemize}
	% segment-id: o014.li2.chapter1.subquotients.i001
	\item Morfisme $f: X \to Y$ disebut \emph{monomorfisme}, dan juga ditulis
	$f: X \hookrightarrow Y$, jika hukum kanselasi kiri berlaku:
	\[ \forall T \in \Obj(\mathcal{C}), \; \forall g,h: T \to X , \quad fg=fh \iff g=h; \]
	% segment-id: o014.li2.chapter1.subquotients.i002
	\item Morfisme $f: X \to Y$ disebut \emph{epimorfisme}, dan juga ditulis
	$f: X \twoheadrightarrow Y$, jika hukum kanselasi kanan berlaku:
	\[ \forall T \in \Obj(\mathcal{C}), \; \forall g,h: Y \to T, \quad gf=hf \iff g=h. \]
\end{itemize}

% segment-id: o014.li2.chapter1.subquotients.x001
\index{monomorfisme (monomorphism)}
\index{epimorfisme (epimorphism)}

% segment-id: o014.li2.chapter1.subquotients.p002
Dengan segera terlihat bahwa $f: X \to Y$ merupakan monomorfisme jika dan
hanya jika $f_*: \Hom(T, X) \to \Hom(T, Y)$ injektif untuk setiap $T$;
$f: X \to Y$ merupakan epimorfisme jika dan hanya jika
$f^*: \Hom(Y, T) \to \Hom(X, T)$ injektif untuk setiap $T$. Sifat
monomorfisme dan epimorfisme saling dual: $f$ merupakan monomorfisme jika dan
hanya jika $f^{\opp}$ merupakan epimorfisme.

% segment-id: o014.li2.chapter1.subquotients.e001
\begin{contoh}
	Jika penyama $\Ker(f,g)$ dari sepasang morfisme $f, g: X \to Y$ ada,
	maka sifat universal menunjukkan bahwa morfisme kanonik
	$\iota: \Ker(f,g) \to X$ merupakan monomorfisme. Secara dual, jika
	kopenyama $\Coker(f,g)$ ada, maka morfisme kanonik
	$p: Y \to \Coker(f,g)$ merupakan epimorfisme. Pengamatan sederhana lainnya
	ialah sebagai berikut: jika $\mathcal{C}$ mempunyai objek nol, yang menurut
	konvensi dinotasikan dengan $0$, maka $0 \to X$ merupakan monomorfisme dan
	$X \to 0$ merupakan epimorfisme. Pembaca dipersilakan memeriksa
	rincian-rincian ini.
\end{contoh}

% segment-id: o014.li2.chapter1.subquotients.pr001
\begin{proposisi}\label{prop:epi-mono-isomorphism}
	Tinjau morfisme-morfisme $X \xrightarrow{a} Y \xrightarrow{b} Z$. Andaikan bahwa
	$b$ merupakan monomorfisme atau $a$ merupakan epimorfisme. Maka
	$ba: X \to Z$ merupakan isomorfisme jika dan hanya jika $a$ dan $b$
	keduanya merupakan isomorfisme.
\end{proposisi}

% segment-id: o014.li2.chapter1.subquotients.pf001
\begin{bukti}
	Salah satu arah sudah jelas. Selanjutnya, andaikan bahwa $ba$ merupakan
	isomorfisme dan $b$ merupakan monomorfisme. Tetapkan
	$c := a(ba)^{-1}: Z \to Y$; maka $bc = \identity_Z$. Selain itu,
	$b(cb) = (bc)b = b$. Karena $b$ merupakan monomorfisme, hukum kanselasi
	kiri memberikan $cb = \identity_Y$. Jadi, $b$ merupakan isomorfisme, dan
	karena itu $a$ juga merupakan isomorfisme. Dengan bekerja di
	$\mathcal{C}^{\opp}$, diperoleh kasus ketika $a$ merupakan epimorfisme.
\end{bukti}

% segment-id: o014.li2.chapter1.subquotients.p003
Dalam kategori yang lazim digunakan dalam aljabar, seperti $\cate{Set}$,
$\cate{Ab}$, dan $R\dcate{Mod}$, sifat monomorfisme dan epimorfisme suatu
morfisme masing-masing ekuivalen dengan sifat injektif dan surjektifnya
sebagai pemetaan. Dalam contoh-contoh ini, kita juga dapat membahas lebih jauh
himpunan bagian, subgrup, dan sebagainya. Konsep tersebut dapat diperluas ke
kategori umum, tetapi diperlukan satu langkah tambahan. Kita akan
merumuskannya dengan bahasa himpunan terurut parsial.

% segment-id: o014.li2.chapter1.subquotients.p004
Tinjau kategori $\mathcal{C}$ dan sebuah objek $X$ di dalamnya. Untuk sepasang
monomorfisme $f_i: S_i \hookrightarrow X$, dengan $i = 1, 2$, jika terdapat
$g: S_1 \to S_2$ sehingga $f_2 g = f_1$, tuliskan
$(S_1, f_1) \subset (S_2, f_2)$, atau cukup $S_1 \subset S_2$. Perhatikan
bahwa sifat monomorfisme $f_2$ mengakibatkan $g$, jika ada, bersifat unik dan
juga merupakan monomorfisme. Relasi $\subset$ yang baru didefinisikan ini
belum merupakan urutan parsial. Jika $S_1 \subset S_2$ dan
$S_2 \subset S_1$, kita katakan bahwa $S_1$ dan $S_2$ ekuivalen. Hal ini juga
setara dengan adanya diagram komutatif

% segment-id: o014.li2.chapter1.subquotients.g001
\[\begin{tikzcd}
	S_1 \arrow[r, "g_1"] \arrow[hookrightarrow, rd, "f_1"'] & S_2 \arrow[hookrightarrow, d, "f_2"] \arrow[r, "g_2"] & S_1 \arrow[hookrightarrow, ld, "f_1"] \\
	& X &
\end{tikzcd}\]

% segment-id: o014.li2.chapter1.subquotients.p005
Keunikan memberikan $g_2 g_1 = \identity_{S_1}$, dan dengan cara yang sama
$g_1 g_2 = \identity_{S_2}$. Dengan kata lain, $S_1$ dan $S_2$ ekuivalen jika
dan hanya jika terdapat isomorfisme $g: S_1 \rightiso S_2$ sehingga
$f_2 g = f_1$; isomorfisme ini pun unik. Tentu saja, inilah argumen baku
dalam aljabar.

% segment-id: o014.li2.chapter1.subquotients.p006
Secara serupa, untuk sepasang epimorfisme
$f_i: X \twoheadrightarrow Q_i$, dengan $i = 1, 2$, jika terdapat
$g: Q_2 \to Q_1$ sehingga $g f_2 = f_1$, maka $g$ unik dan juga merupakan
epimorfisme; tuliskan $Q_1 \twoheadleftarrow Q_2$. Jika
$Q_1 \twoheadleftarrow Q_2$ dan $Q_2 \twoheadleftarrow Q_1$, kita katakan bahwa
$Q_1$ dan $Q_2$ ekuivalen. Hal ini juga setara dengan adanya isomorfisme unik
$g: Q_2 \rightiso Q_1$ sehingga $g f_2 = f_1$.

% segment-id: o014.li2.chapter1.subquotients.p007
Relasi biner $\subset$ (atau $\twoheadleftarrow$) menginduksi urutan parsial
pada kelas-kelas ekuivalensi. Di sini kita meminjam simbol inklusi himpunan
$\subset$. Meskipun hal ini dapat menimbulkan kerancuan, notasi tersebut
selaras dengan notasi bagi substruktur lain dalam aljabar, seperti subgrup,
submodul, dan subruang.

% segment-id: o014.li2.chapter1.subquotients.d001
\begin{definisi}\label{def:sub-quot-obj}
	\index{subobjek (subobject)} \index[sym1]{SubX@$\mathrm{Sub}_X$}
	\index{objek kuosien (quotient object)} \index[sym1]{QuotX@$\mathrm{Quot}_X$}
	Tetapkan kategori $\mathcal{C}$ dan sebuah objek $X$ di dalamnya. Kelas
	ekuivalensi dari monomorfisme berbentuk $S \hookrightarrow X$ (atau dari epimorfisme
	berbentuk $X \twoheadrightarrow Q$) menurut relasi ekuivalensi di atas
	disebut \emph{subobjek} (atau \emph{objek kuosien}) dari $X$. Kelas-kelas
	itu membentuk himpunan terurut parsial $(\mathrm{Sub}_X, \subset)$ (atau
	$(\mathrm{Quot}_X, \twoheadleftarrow)$), dengan
	$\identity_X: X \to X$ sebagai unsur maksimal yang tunggal.
\end{definisi}

% segment-id: o014.li2.chapter1.subquotients.p008
Subobjek dan objek kuosien saling dual: $(\mathrm{Sub}_X, \subset)$ yang
didefinisikan untuk $\mathcal{C}$ sama dengan
$(\mathrm{Quot}_X, \twoheadleftarrow)$ yang didefinisikan untuk
$\mathcal{C}^{\opp}$. Demi kemudahan pemaparan, bagian ini terutama membahas
subobjek. Karena isomorfisme $g$ dalam relasi ekuivalensi bersifat unik, dalam
argumen kita dapat dengan aman memilih wakil $f: S \hookrightarrow X$ dari
suatu kelas ekuivalensi, atau bahkan menuliskannya cukup sebagai $S$.
Ingatlah bahwa $f_*: \Hom(\cdot, S) \to \Hom(\cdot, X)$ bersifat injektif.

% segment-id: o014.li2.chapter1.subquotients.lm001
\begin{lema}\label{prop:subobject-order}
	Untuk setiap dua subobjek $S_1$ dan $S_2$ dari $X$, berlaku
	\begin{equation*}
		S_1 \subset S_2 \iff \forall T \in \Obj(\mathcal{C}), \; \Hom(T, S_1) \subset \Hom(T, S_2),
	\end{equation*}
	dengan $\Hom(T, S_1)$ dan $\Hom(T, S_2)$ pada ruas kanan dipahami sebagai
	himpunan bagian dari $\Hom(T, X)$.
\end{lema}

% segment-id: o014.li2.chapter1.subquotients.pf002
\begin{bukti}
	Arah $\implies$ mudah. Untuk arah $\impliedby$, ambil $T := S_1$ dan
	perhatikan $\identity_{S_1} \in \Hom(S_1, S_1)$. Di bawah penyertaan
	$\Hom(S_1,S_1) \subset \Hom(S_1,X)$, unsur ini bersesuaian dengan $f_1$.
	Hipotesis lalu memberikan $g \in \Hom(S_1,S_2)$ yang citranya dalam
	$\Hom(S_1,X)$ adalah $f_1$, yakni $f_2g=f_1$. Inilah morfisme yang dicari.
\end{bukti}

% segment-id: o014.li2.chapter1.subquotients.p009
Andaikan bahwa $\mathcal{C}$ mempunyai objek nol dan
$(X_i)_{i \in I}$ merupakan suatu keluarga objek. Untuk setiap
$(i,j) \in I \times I$, definisikan $\delta_{ij} \in \Hom(X_i, X_j)$ sebagai
berikut: jika $i=j$, ambil $\delta_{ij} = \identity$; jika tidak, tetapkan
$\delta_{ij}$ sebagai morfisme nol. Jika produk dan koproduk dari
$(X_i)_{i \in I}$ ada, maka $(\delta_{ij})_{i,j}$ menentukan morfisme kanonik

% segment-id: o014.li2.chapter1.subquotients.q001
\begin{equation}\label{eqn:coprod-prod-delta}
	\delta: \coprod_{i \in I} X_i \to \prod_{i \in I} X_i.
\end{equation}

% segment-id: o014.li2.chapter1.subquotients.p010
Dalam kategori aditif, yang akan segera kita tinjau kembali, $\delta$ selalu
merupakan isomorfisme apabila $I$ hingga. Hal ini belum tentu berlaku dalam
kategori umum.

% segment-id: o014.li2.chapter1.subquotients.lm002
\begin{lema}\label{prop:mono-pullback}
	Andaikan $X \to Z$ merupakan monomorfisme. Tarik baliknya sepanjang
	$Y \to Z$, yaitu $X \dtimes{Z} Y \to Y$ (andaikan ada), juga merupakan
	monomorfisme. Secara dual, dorong keluar suatu epimorfisme juga merupakan
	epimorfisme.
\end{lema}

% segment-id: o014.li2.chapter1.subquotients.pf003
\begin{bukti}
	Berdasarkan dualitas, cukup menangani kasus monomorfisme. Misalkan
	$p_2: X \dtimes{Z} Y \to Y$ adalah morfisme kanonik yang diberikan oleh
	produk serat. Untuk setiap objek $T$, masalahnya cukup direduksi menjadi
	pembuktian bahwa
	% segment-id: o014.li2.chapter1.subquotients.g002
	\[\begin{tikzcd}
		\Hom(T,X) \dtimes{\Hom(T,Z)} \Hom(T,Y) & \\
		\Hom\left(T, X \dtimes{Z} Y \right) \arrow[leftrightarrow, u, "1:1", "\text{sifat universal}"'] \arrow[r, "{(p_2)_*}"'] & \Hom(T, Y)
	\end{tikzcd} \]
	bersifat injektif. Akan tetapi, karena $X \to Z$ merupakan monomorfisme,
	pemetaan $\Hom(T, X) \to \Hom(T, Z)$ bersifat injektif, sedangkan pemetaan
	dari kiri atas ke kanan bawah hanyalah proyeksi alami.
\end{bukti}

% segment-id: o014.li2.chapter1.subquotients.lm003
\begin{lema}\label{prop:mono-product}
	Andaikan $\left( f_i: X_i \to Z \right)_{i \in I}$ merupakan suatu keluarga
	monomorfisme. Maka morfisme kanonik dari
	$\prod_{i \in I} (X_i \to Z)$ (andaikan ada) ke $Z$ juga merupakan
	monomorfisme. Secara dual, untuk suatu keluarga epimorfisme
	$\left( g_i: Z \to X_i \right)_{i \in I}$, morfisme kanonik dari $Z$ ke
	$\coprod_{i \in I} (Z \to X_i)$ (andaikan ada) juga merupakan
	epimorfisme.
\end{lema}

% segment-id: o014.li2.chapter1.subquotients.pf004
\begin{bukti}
	Seperti pada Lema~\ref{prop:mono-pullback}, cukup dibuktikan bahwa baris
	kedua dalam diagram
	% segment-id: o014.li2.chapter1.subquotients.g003
	\[\begin{tikzcd}
		\prod_{i \in I} \left( \Hom(T, X_i) \to \Hom(T, Z) \right) & \\
		\Hom\left(T, \prod_{i \in I} (X_i \to Z) \right) \arrow[leftrightarrow, u, "1:1", "\text{sifat universal}"'] \arrow[r] & \Hom(T, Z)
	\end{tikzcd}\]
	bersifat injektif. Hal ini mengikuti dari hipotesis, sebab setiap
	$\Hom(T, X_i) \to \Hom(T, Z)$ bersifat injektif.
\end{bukti}

% segment-id: o014.li2.chapter1.subquotients.p011
Kita kembali ke pembahasan subobjek dan objek kuosien.

% segment-id: o014.li2.chapter1.subquotients.d002
\begin{definisi}\label{def:subquotient}
	\index{subkuosien (subquotient)}
	Diberikan sembarang kategori $\mathcal{C}$ beserta objek-objek $X$ dan $Y$
	di dalamnya. Jika $Y$ dapat direalisasikan sebagai objek kuosien dari suatu
	subobjek dari $X$, maka $Y$ disebut \emph{subkuosien} dari $X$.
\end{definisi}

% segment-id: o014.li2.chapter1.subquotients.p012
Jika yang ditinjau adalah subobjek dari objek kuosien, hal itu sama dengan
membahas subkuosien dalam $\mathcal{C}^{\opp}$. Hasil berikut menunjukkan
bahwa, untuk suatu kelas kategori, kedua cara tersebut memberikan pengertian
yang sama. Kelas ini mencakup kategori abelian yang akan dipelajari kemudian;
lihat Korolari~\sourcecrossref{prop:Abel-cat-subquotient}{2.1.7}.

% segment-id: o014.li2.chapter1.subquotients.pr002
\begin{proposisi}\label{prop:subquotient-equiv}
	Andaikan kategori $\mathcal{C}$ mempunyai semua produk serat hingga dan
	koproduk serat hingga, tarik balik epimorfisme tetap merupakan
	epimorfisme, dan dorong keluar monomorfisme tetap merupakan monomorfisme.
	Maka subkuosien dapat didefinisikan secara ekuivalen sebagai subobjek dari
	objek kuosien.
\end{proposisi}

% segment-id: o014.li2.chapter1.subquotients.pf005
\begin{bukti}
	Tinjau diagram subkuosien $Y \twoheadleftarrow Z \hookrightarrow X$ dalam
	$\mathcal{C}$. Tetapkan $W := X \dsqcup{Z} Y$. Morfisme $X \to W$ adalah
	dorong keluar dari $Z \to Y$, sehingga merupakan epimorfisme menurut
	Lema~\ref{prop:mono-pullback}; morfisme $Y \to W$ adalah dorong keluar dari
	$Z \to X$, sehingga merupakan monomorfisme menurut hipotesis. Jadi,
	$Y \hookrightarrow W \twoheadleftarrow X$, yang memberikan subkuosien dari
	$X$ dalam $\mathcal{C}^{\opp}$. Karena syarat-syarat tersebut swa-dual,
	subkuosien dalam $\mathcal{C}^{\opp}$ juga memberikan subkuosien dalam
	$\mathcal{C}$.
\end{bukti}

% segment-id: o014.li2.chapter1.subquotients.p013
Setelah sifat ``objek kuosien dari subobjek = subobjek dari objek kuosien'' di
atas tersedia, dapat disimpulkan bahwa subkuosien dari subkuosien tetap
merupakan subkuosien.

% segment-id: o014.li2.chapter1.subquotients.p014
Definisi subkuosien tidak menentukan bagaimana subkuosien itu direalisasikan
sebagai objek kuosien dari suatu subobjek. Misalnya, untuk
$\mathcal{C} = \cate{Ab}$, $\Z/2\Z$ merupakan subkuosien dari $\Z/4\Z$:
objek tersebut dapat direalisasikan sebagai objek kuosien dari $\Z/4\Z$,
tetapi juga sebagai subobjek $2\Z/4\Z$ dari $\Z/4\Z$.
